\IfFileExists{papers/formalized-vs-literature-preamble.tex}{\documentclass[11pt]{article}
\usepackage[utf8]{inputenc}
\usepackage[T1]{fontenc}
\usepackage{amsmath,amssymb,mathtools}
\usepackage{booktabs,tabularx,array}
\usepackage{enumitem}
\usepackage[margin=1in]{geometry}
\usepackage{xcolor}
\usepackage[hidelinks]{hyperref}
\usepackage{microtype}
\setlength{\parindent}{0pt}
\setlength{\parskip}{0.55em}
\newcommand{\Lean}[1]{\texttt{\nolinkurl{#1}}}
\newcommand{\RepoFile}[1]{\texttt{\nolinkurl{#1}}}
\newcommand{\Class}[1]{\textbf{#1}}
\newcommand{\Top}{\mathord{\top}}
\newcommand{\R}{\mathbb{R}}
\newcommand{\ENN}{\mathbb{R}_{\ge 0}^{\infty}}
\newcommand{\KL}{\operatorname{KL}}
\newcommand{\W}{\mathsf{W}}
\newcommand{\statusline}[2]{%
  \begin{center}\fbox{\begin{minipage}{0.92\linewidth}\textbf{Completion class:} #1\\[2pt]
  \textbf{Strongest caveat:} #2\end{minipage}}\end{center}}
\newenvironment{comparison}{\tabularx{\linewidth}{>{\bfseries}p{0.25\linewidth}X}}{\endtabularx}
\newcommand{\snapshot}{\texttt{902f51aa01a737af68f6feb71c83263cc9b1f4d9}}
}{\documentclass[11pt]{article}
\usepackage[utf8]{inputenc}
\usepackage[T1]{fontenc}
\usepackage{amsmath,amssymb,mathtools}
\usepackage{booktabs,tabularx,array}
\usepackage{enumitem}
\usepackage[margin=1in]{geometry}
\usepackage{xcolor}
\usepackage[hidelinks]{hyperref}
\usepackage{microtype}
\setlength{\parindent}{0pt}
\setlength{\parskip}{0.55em}
\newcommand{\Lean}[1]{\texttt{\nolinkurl{#1}}}
\newcommand{\RepoFile}[1]{\texttt{\nolinkurl{#1}}}
\newcommand{\Class}[1]{\textbf{#1}}
\newcommand{\Top}{\mathord{\top}}
\newcommand{\R}{\mathbb{R}}
\newcommand{\ENN}{\mathbb{R}_{\ge 0}^{\infty}}
\newcommand{\KL}{\operatorname{KL}}
\newcommand{\W}{\mathsf{W}}
\newcommand{\statusline}[2]{%
  \begin{center}\fbox{\begin{minipage}{0.92\linewidth}\textbf{Completion class:} #1\\[2pt]
  \textbf{Strongest caveat:} #2\end{minipage}}\end{center}}
\newenvironment{comparison}{\tabularx{\linewidth}{>{\bfseries}p{0.25\linewidth}X}}{\endtabularx}
\newcommand{\snapshot}{\texttt{902f51aa01a737af68f6feb71c83263cc9b1f4d9}}
}
\title{Wasserstein DRSB card bound: formalized result versus the literature}
\author{}
\date{Formalization snapshot \snapshot}
\begin{document}
\maketitle
\statusline{\Class{S/G}: complete quadratic weak-duality specialization built from a reusable general-cost kernel}{This is a one-sided bound for one feasible law, not Gao--Kleywegt or Blanchet--Murthy strong duality, and the radius is a budget on squared transport cost rather than the unsquared $W_2$ radius.}
\section{Comparison target}
The production declaration is
\[
  \Lean{Drsb.wdrsb_cost_bound}
  \quad\text{in}\quad
  \RepoFile{Drsb/Basic.lean}.
\]
It supports the evaluation-card claim that a particular perturbed source law inside the quadratic transport ambiguity set has expected value no larger than the scalar Wasserstein-DRO dual bound.

The closest literature sources are Blanchet--Murthy (2019), Theorem~1 and Remark~1, and Gao--Kleywegt (2023), Proposition~1 and Theorem~1.  The local reconstructions are
\RepoFile{prose/distilled_literature/BlanchetMurthy2019_theorem1.tex} and
\RepoFile{prose/distilled_literature/GaoKleywegt2023_theorem1.tex}.

\section{Literature statement used by the card}
For a nominal law $p_0$, reward $V$, nonnegative cost $c$, and budget $\delta$, weak duality gives, for every feasible $\mu$,
\[
  \int V\,d\mu
  \le
  \inf_{\lambda\ge 0}
  \left\{\lambda\delta+
  \int \sup_x\bigl(V(x)-\lambda c(x,y)\bigr)\,p_0(dy)\right\}.
\]
The card uses $c(x,y)=\lVert x-y\rVert^2$.  If a paper writes $W_2(\mu,p_0)\le r$, then the corresponding Lean budget is $\delta=r^2$.

\section{What Lean proves}
The theorem assumes:
\begin{itemize}[leftmargin=2em]
\item $\mu$ belongs to the real-valued quadratic-cost ball \Lean{wassersteinBall p0 epsilon};
\item $V$ is integrable under $\mu$;
\item every pointwise dual envelope is bounded above and integrable under $p_0$;
\item both $\mu$ and $p_0$ have finite second moments.
\end{itemize}
It then proves the displayed one-sided bound.  No optimal transport plan is assumed to attain the infimum.  Instead, the proof selects plans with cost below $\epsilon+\eta$ and sends $\eta\downarrow0$.  The second-moment assumptions imply integrability of the quadratic cost for every coupling.

The load-bearing paper-agnostic theorem is
\Lean{ForMathlib.OT.expect_le_dualIntegrand_add_lam_couplingCost} in
\RepoFile{ForMathlib/OptimalTransport/WeakDuality.lean}.

\section{Axis-by-axis comparison}
\begin{comparison}
Source scope & General nonnegative lower-semicontinuous or metric-power costs, depending on the source theorem.\\
Lean scope & Quadratic norm cost in the card wrapper; the underlying Lagrangian inequality is general-cost.\\
Value type & The card uses real Bochner integrals, real $sInf$, and real $sSup$.\\
Radius convention & $\epsilon$ bounds $W_2^2$, not $W_2$.\\
Duality level & One feasible-law inequality only.\\
Attainment & Neither primal nor dual attainment is claimed or needed.\\
Continuum DRSB dependency & None: $V$ is an abstract function.\\
\end{comparison}

\section{Why the second-moment hypotheses remain}
The current real-valued \Lean{couplingCost} uses a Bochner integral.  A nonintegrable cost can therefore collapse to the real integral's junk value.  Finite moments prevent this pathology for the quadratic cost.  This is a formalization artifact, not a hypothesis in the natural extended-valued formulation.  A canonical $\ENN$ transport layer should eventually remove it.

\section{Explicit nonclaims}
This result does not prove equality of primal and dual values, finiteness of either optimization problem, existence of a worst-case distribution, or construction of the stochastic-control value function $V$.  Those are separate results and campaigns.
\end{document}
